%! Author = chtompki
%! Date = 1/14/26
\documentclass[12pt]{amsart}
% Default margins are too wide all the way around. I reset them here
\setlength{\hoffset}{-.75in}
\setlength{\textwidth}{6.5in}
\setlength{\voffset}{-.5in}
\setlength{\textheight}{9.0in}
\usepackage{amsmath}
\usepackage{leftindex}
\usepackage{amssymb}
\usepackage{amsthm}
\usepackage{enumitem}
\usepackage{setspace}

\newenvironment{solution}
{\begin{proof}[Solution]}
{\end{proof}}

\title{Combinatorics\\Assignment 5}
\author{Rob Tompkins}
\date{\today}

\begin{document}
    \maketitle
    \date{\today}
    \begin{onehalfspace}
        \section{Problem 1}
        In how many ways can we build a 7-beaded circular necklace if there are 5 available types of gemstones
        (repetition) allowed and all rotations  and reflections of a given necklace are equivalent.
        What if we have $q$ gemstones?

        \begin{solution}
            Our goal is an application of Burnside's Lemma here.
            Now, for Burnside's Lemma, we need to start
            with a finite group $G$ and a set $X$ and $G$ acting on $X$.
            For each $g\in G$, set $X^g$ to be the elements of $X$ that are fixed by $g$'s action on $X$.
            That is to say we set $X^g = \{x\in X: g\cdot x = x\}$.
            Burnside's Lemma then asserts that the number of $G$ orbits of $X$, denoted $|G\\X|$ is given by
            \[
                |G\\X| = \frac{1}{|G|}\sum_{g\in G} |X^g|.
            \]

            We start by noticing that in our problem, we are dealing with the group of rotations and reflections
            on the 7-cycle, which is commonly referred to as the dihedral group on 7 elements, $D_7$.
            Notice that the size of $|D_7|=14$.

            Assume we have $q$ gemstones for the sake of generality.
            Now, we must account for all of the elements of $D_7$ in order to account for the summation in Burnside's
            Lemma.

            So we begin with the identity rotation which fixes our 7 places with $q$ colorings counting $q^7$ occurrences.
            We then proceed to counting the non-trivial rotations of which there are 6, so we get another $6q$ coloring
            occurrences there.
            Lastly, we consider the reflections, for which there are 7 axis of symmetry each bead on the axis can be
            any of $q$ colors and the remaining 6 beads are set up in 3 pairs where each of the 3 pairs must match in color.
            Thus, we count this number of colorings as $7\cdot q\cdot q^3=7q^4$.

            So when we put all this together we see that
            \[
                |D_7\\[q]| = \frac{1}{14}(q^7 + 6q + 7q^4).
            \]
            So for 5 colors we have]
            \[
                |D_7\\[5]| = \frac{1}{14}(5^7 + 6\cdot 5 + 7\cdot 5^4) = 5895
            \]


        \end{solution}

        \section{Problem 2}
        A rectangular prism has equilateral triangles for its base and top, and three rectangular lateral faces.
        In how many ways can we color if the faces can be colored using any one of the 5 colors.

        \begin{solution}

        \end{solution}

        \section{Problem 3}
        Find the pattern inventory of the red-blue vertex colorings of a tetrahedron.
        Repeat the problem with red-blue edge colorings.

        \begin{solution}

        \end{solution}

        \section{Problem 4}
        How many 10 red, 12 blue, 3 green squares are there in the coloring of a $5\times 5$ chessboard?

        \begin{solution}

        \end{solution}

        \section{Problem 5}
        Find the pattern inventory of the red-blue vertex colorings of a floating cube (i.e. the cube can freely move
        in space).

        \begin{solution}

        \end{solution}
    \end{onehalfspace}
\end{document}
